\chapter{Couche 1 — L'application web}

Six briques pour transformer du code Python en un site que des gens peuvent
utiliser. Chacune règle un problème que la précédente laissait ouvert.

\begin{figure}[htbp]
\centering
\resizebox{\textwidth}{!}{%
\begin{tikzpicture}[font=\scriptsize,
  b/.style={draw=grisligne,fill=grisfond,rounded corners=2pt,minimum height=13mm,
            minimum width=2.3cm,align=center},
  f/.style={-{Stealth[length=2mm]},gristexte}]
\node[b] (n) at (0,0) {Navigateur};
\node[b] (x) at (3.2,0) {\textbf{nginx}\\\textcolor{gristexte}{TLS, statiques}};
\node[b] (g) at (6.4,0) {\textbf{Gunicorn}\\\textcolor{gristexte}{4 ouvriers}};
\node[b] (fl) at (9.6,0) {\textbf{Flask}\\\textcolor{gristexte}{routes}};
\node[b] (j) at (12.8,0) {\textbf{Jinja2}\\\textcolor{gristexte}{gabarits}};
\draw[f] (n) -- node[above]{HTTPS} (x);
\draw[f] (x) -- node[above]{HTTP local} (g);
\draw[f] (g) -- node[above]{WSGI} (fl);
\draw[f] (fl) -- node[above]{rendu} (j);
\end{tikzpicture}}
\caption{Le chemin d'une requête. Chaque maillon existe parce que le suivant ne
sait pas faire son travail.}
\end{figure}

% -----------------------------------------------------------------------------
\outil{Flask}{brique 4 · le cadre web}

\quoi
Un \emph{micro}-framework : il fournit le routage (associer une URL à une
fonction), les gabarits et les sessions. Il ne fournit ni base de données, ni
administration, ni authentification --- on choisit soi-même.

\pourquoi
L'autre grand choix Python est \textbf{Django}, qui fournit tout : ORM,
interface d'administration, comptes utilisateurs, formulaires. Pour un projet
d'entreprise classique, Django va plus vite.

Ici, c'était un mauvais choix pour une raison précise : le sujet demande de
comprendre et de montrer chaque mécanisme. Avec Django, l'authentification est
une boîte noire qu'on active ; avec Flask, on écrit le hachage, la session, le
verrouillage après cinq échecs, et on peut les expliquer. Le projet est un
exercice d'ingénierie, pas une course à la livraison.

Deuxième raison : Django impose son ORM. Nous voulions pouvoir piloter la base
finement --- notamment l'index unique partiel qui garantit l'unicité des ports.

\role
\fichier{app.py} porte les routes HTML et les sessions ;
\fichier{api.py} porte les routes JSON. Une particularité : \textbf{les mêmes
URL servent les deux}, selon l'en-tête \cmd{Accept} de l'appelant. Il n'y a pas
un site d'un côté et une API de l'autre, mais une logique métier unique avec
deux rendus.

\begin{lstlisting}[language=Python,caption={\fichier{2\_webapp/app.py}}]
def veut_json() -> bool:
    # Un client qui accepte TOUT (Accept: */*, cas de curl sans en-tete)
    # n'est pas un client d'API : seul celui qui demande du JSON sans accepter
    # le HTML est traite comme tel.
    accepte = request.accept_mimetypes
    return accepte.accept_json and not accepte.accept_html
\end{lstlisting}

\apport
La plateforme devient atteignable depuis un navigateur. Et, grâce à la
négociation de contenu, elle est simultanément pilotable par \cmd{curl} sans
une ligne de code en plus --- ce qui satisfait l'exigence d'API du sujet sans
dédoubler le travail.

% -----------------------------------------------------------------------------
\outil{Jinja2}{brique 5 · les gabarits}

\quoi
Le moteur qui remplit des fichiers HTML avec des variables. Dans le gabarit on
écrit \verb|{{ machine.nom }}| et \verb|{% for m in machines %}| ; Jinja2
produit le HTML final.

\pourquoi
L'alternative serait de construire le HTML par concaténation de chaînes en
Python. Outre l'illisibilité, c'est une faille : une variable contenant
\cmd{<script>} serait injectée telle quelle dans la page. Jinja2
\textbf{échappe automatiquement} tout ce qu'il insère --- c'est une protection
contre les injections de script obtenue sans y penser.

Il vient avec Flask, donc aucun choix à faire ; mais le mentionner compte,
parce que l'échappement automatique est une des raisons pour lesquelles le
projet n'a aucune faille XSS.

\role
Les gabarits de \fichier{2\_webapp/templates/} : la page de connexion,
l'inscription, et le tableau de bord qui fait à lui seul 358 lignes. Les
conditions y servent à n'afficher une section que si elle a du contenu ---
« Archives de mon travail » n'apparaît que s'il y a des archives.

\apport
Des pages lisibles à écrire, et une classe entière d'attaques écartée par
défaut.

% -----------------------------------------------------------------------------
\outil{Werkzeug}{brique 6 · la plomberie HTTP}

\quoi
La bibliothèque sur laquelle Flask est construit. Elle gère le détail du
protocole HTTP, et fournit deux outils qu'on utilise directement :
le hachage des mots de passe et le \emph{middleware} \cmd{ProxyFix}.

\pourquoi
On ne l'a pas choisie : elle arrive avec Flask. Mais deux de ses fonctions sont
employées délibérément plutôt que réécrites.

\cmd{generate\_password\_hash} applique \textbf{pbkdf2-sha256} avec un sel
aléatoire et des milliers d'itérations. Écrire soi-même un hachage de mot de
passe est l'erreur classique : on choisit un algorithme trop rapide, on oublie
le sel, et la base devient exploitable en cas de fuite.

\cmd{ProxyFix} règle un problème qu'on ne voit qu'en production : derrière
nginx, toutes les requêtes semblent venir de \cmd{127.0.0.1}. Sans lui, notre
verrouillage anti-force-brute compterait toutes les tentatives comme venant de
la même adresse --- et bloquerait tout le monde au cinquième échec de n'importe
qui.

\role
Hachage à l'inscription et vérification à la connexion ; \cmd{ProxyFix} activé
dans \fichier{app.py} pour que l'adresse réelle du client remonte.

\apport
Des mots de passe stockés correctement, et un anti-force-brute qui vise la
bonne personne.

% -----------------------------------------------------------------------------
\outil{Gunicorn}{brique 7 · le serveur d'application}

\quoi
Un serveur WSGI : il reçoit les requêtes HTTP et les passe à l'application
Python. Il lance plusieurs processus ouvriers qui se partagent le travail.

\pourquoi
Flask embarque un serveur de développement, mais il traite \textbf{une requête
à la fois} et affiche lui-même un avertissement : ne pas utiliser en
production. Or créer un conteneur prend plusieurs secondes ; avec un seul
processus, un utilisateur qui loue une machine bloquerait tous les autres.

Les alternatives sont uWSGI (plus puissant, bien plus complexe à configurer) et
Waitress (plus simple, mais mono-processus par défaut). Gunicorn est le milieu
raisonnable : une ligne de commande suffit, et il redémarre les ouvriers qui
plantent.

\role
Quatre ouvriers, lancés par l'unité systemd du contrôleur :

\begin{lstlisting}[language=bash]
gunicorn --preload -w 4 -b 127.0.0.1:5000 app:app
\end{lstlisting}

\begin{piege}{Quatre processus qui démarrent en même temps}
Chacun importe l'application, qui crée les tables si elles manquent. Les quatre
le font simultanément, et parfois l'un échoue avec \cmd{table users already
exists}. Sur un poste mono-processus, le bug n'apparaît jamais.

Deux correctifs : la création tolère désormais l'erreur « existe déjà », et
l'option \cmd{--preload} fait importer l'application \emph{une fois} avant de
dupliquer le processus. Nous avons reproduit la panne avec huit processus
simultanés --- six sur huit démarraient ; après correction, huit sur huit.
\end{piege}

\apport
La plateforme supporte plusieurs utilisateurs simultanés. Avant Gunicorn, une
location en cours gelait le site pour tout le monde.

% -----------------------------------------------------------------------------
\outil{nginx}{brique 8 · le proxy inverse}

\quoi
Un serveur web placé \emph{devant} l'application. Le navigateur lui parle ; il
relaie vers Gunicorn.

\pourquoi
Gunicorn sait répondre à du HTTP, mais ne sait pas faire trois choses
essentielles : terminer le TLS, servir efficacement les fichiers statiques, et
absorber les requêtes malformées.

L'alternative est Apache, qui sait tout faire aussi mais consomme davantage et
se configure plus lourdement. Caddy aurait apporté les certificats
automatiques, ce qui n'a pas de sens sur un réseau privé sans nom de domaine
public.

\role
Le rôle Ansible \cmd{reverse\_proxy} installe nginx avec :

\begin{itemize}
\item \textbf{TLS 1.3} et le certificat du nœud ;
\item les fichiers statiques servis directement, sans réveiller Python ;
\item les en-têtes de sécurité : \cmd{Strict-Transport-Security},
  \cmd{X-Content-Type-Options}, \cmd{X-Frame-Options} ;
\item \cmd{X-Forwarded-For}, que \cmd{ProxyFix} relira côté Flask.
\end{itemize}

\apport
Le site passe en HTTPS, les images et le CSS ne coûtent plus de temps Python, et
l'application n'est plus exposée directement au réseau.

% -----------------------------------------------------------------------------
\outil{Tailwind CSS, précompilé}{brique 9 · l'apparence}

\quoi
Une bibliothèque de style où l'on compose l'apparence avec de petites classes
dans le HTML, plutôt qu'en écrivant du CSS à part.

\pourquoi
L'usage courant est de charger Tailwind depuis un CDN : une ligne, et ça marche.
C'était impossible ici.

La politique de sécurité de contenu du projet est \cmd{default-src 'self'} :
\textbf{aucune ressource externe} n'a le droit d'être chargée. Un CDN est
précisément une ressource externe. Il a donc fallu \textbf{précompiler} Tailwind
en un fichier CSS servi par nginx, et héberger les polices localement.

\role
\fichier{2\_webapp/static/css/tailwind.css} est le résultat de la compilation ;
\fichier{tailwind.src.css} en est la source. Les polices IBM Plex Sans sont dans
\fichier{static/fonts/}.

\apport
Une interface soignée \emph{sans} affaiblir la politique de sécurité. C'est le
genre de compromis où le choix facile (le CDN) aurait coûté une faille entière :
avec un CDN autorisé, une injection de script redevient exploitable.

\gros[0.55]{zoom-connexion.png}{La page de connexion. Tout le style vient d'un
fichier CSS local : la politique de sécurité interdit toute ressource externe.}

% -----------------------------------------------------------------------------
\section*{Ce que la plateforme sait faire après la couche 1}

Elle a un site en HTTPS, capable de servir plusieurs utilisateurs, avec des
comptes et des sessions. Mais elle ne retient rien : à chaque redémarrage, tout
disparaît. C'est l'objet de la couche suivante.
